\section{Discussion}
\label{sec:discussion}

OriginBlame has four limitations. (1)~Incremental adoption is difficult: users cannot retroactively add provenance to existing datasets. (2)~The parser ecosystem is immature: only a MediaWiki parser is currently available, though our cross-domain evaluation (\S\ref{sec:cross-domain}) confirms that the core provenance tracking is domain-agnostic and works with arbitrary attribution sources like git blame. (3)~Single-hop provenance: the system tracks only direct mappings from raw data to final output. (4)~The index must be rebuilt when data changes; indexed variants do not consistently outperform the full-scan path due to rayon parallelism amortizing the scan cost.

Future work will pursue three directions: expanding the parser ecosystem to support Common Crawl archives, PDF documents, and mainstream NLP dataset formats; exploring incremental adoption paths that allow users to partially embed OriginBlame in existing pipelines; and extending the token-index layer to support per-token attribution (rather than per-document token-count attribution) for integration with token-level unlearning methods.

\paragraph{Privacy and dual-use.}
The authors layer stores contributor names and emails alongside their provenance records. In production deployments handling sensitive contributor data, operators can strip PII from \texttt{.ob/} by removing name and email fields from author records---leaving only the SHA-256 identifier---and maintaining a separate encrypted mapping outside the provenance directory. The operation log (\texttt{.ob/log}) provides an audit trail for revocation and query commands, supporting access-control policies that restrict provenance queries to authorized compliance personnel. We note that fine-grained author attribution could be misused to target individual contributors; deploying organizations should restrict provenance query access accordingly.
